\chapter{CƠ SỞ LÝ THUYẾT VÀ PHƯƠNG PHÁP NGHIÊN CỨU}
\label{ch:phuong-phap}

Chương này trình bày cơ sở lý thuyết, dữ liệu, kiến trúc hệ thống và giao thức đánh giá
của đề tài. Quy trình được thiết kế để kết nối hai nguồn thông tin có bản chất khác nhau:
tin tức tài chính tiếng Việt mô tả sự kiện và kỳ vọng, còn dữ liệu thị trường ghi nhận
kết quả giao dịch theo thời gian. Vì vậy, các bước thu thập, làm sạch, gắn mã cổ phiếu,
ánh xạ thời điểm, tạo nhãn và chia tập được mô tả trước khi trình bày phần huấn luyện
mô hình.

Phương pháp được xây dựng theo hai tầng. Tầng bắt buộc khóa mức sàn gồm đối chứng cơ sở,
mô hình chỉ dùng giá, mô hình hai nhánh, phép loại bỏ đặc trưng và đánh giá cuốn chiếu
theo thời gian. Tầng mở rộng chỉ được xem xét sau đó, gồm các mô hình phức tạp hơn như
Bộ biến đổi hợp nhất theo thời gian, hợp nhất có cổng hoặc biểu diễn sự kiện. Cách tổ
chức này cho phép trả lời từng câu hỏi riêng: dữ liệu có đủ và hợp lệ hay không, mô hình
giá có học được tín hiệu hay không, và cuối cùng cảm xúc có bổ sung thông tin hay không.
Khi một tầng trước chưa được kiểm tra, điểm số của tầng sau không được dùng để che lấp
vấn đề ở tầng dữ liệu.

\section{Tổng quan kiến trúc hệ thống}
\label{sec:pp-kien-truc}

Hệ thống gồm hai nhánh xử lý riêng trước khi hợp nhất. Nhánh cảm xúc nhận tiêu đề và
nội dung tin tức tiếng Việt trong phạm vi được phép sử dụng, tinh chỉnh PhoBERT để sinh
xác suất của ba lớp tích cực, trung tính và tiêu cực, rồi tổng hợp tín hiệu theo mã và
ngày giao dịch. Nhánh giá nhận dữ liệu OHLCV điều chỉnh, tính lợi suất, độ biến động,
khối lượng và các chỉ báo kỹ thuật theo lịch phiên. Sau khi áp dụng quy tắc mốc cắt,
hai nhánh được căn chỉnh bằng khóa gồm mã cổ phiếu và ngày giao dịch.

Mỗi dòng của bảng đặc trưng phải kèm trạng thái dữ liệu cần thiết, chẳng hạn ngày đó có
tin hay không, tin có mốc thời gian hợp lệ hay không và giá có đủ lịch sử hay không.
Bảng đặc trưng thu được dùng để dự báo lớp UP, FLAT hoặc DOWN của phiên kế tiếp. Giá
của phiên kế tiếp chỉ xuất hiện ở cột nhãn sau khi dữ liệu đầu vào đã được khóa, và
không được quay ngược lại để tạo đặc trưng cho ngày quan sát.

\subsection{Nguyên tắc thiết kế}

Kiến trúc được xây dựng theo năm nguyên tắc. Thứ nhất, mỗi bản ghi phải có mốc thời
gian, nguồn gốc và trạng thái ánh xạ để kiểm tra khả năng sử dụng tại thời điểm dự báo.
Thứ hai, tín hiệu cảm xúc được lưu dưới dạng xác suất thay vì chỉ một nhãn cứng, nhằm
giữ lại thông tin về mức độ không chắc chắn và cho phép tính các đại lượng tổng hợp
khác nhau. Thứ ba, mô hình có cảm xúc phải được đặt cạnh mô hình chỉ dùng giá trong
cùng điều kiện, nếu không sẽ không thể quy phần chênh lệch cho thông tin văn bản. Thứ
tư, mọi phép biến đổi có học tham số phải được ước lượng trong cửa sổ huấn luyện tương
ứng. Thứ năm, mô hình mở rộng chỉ được thực hiện sau khi mức sàn đã vượt qua các kiểm
tra về dữ liệu, rò rỉ và độ ổn định. Năm nguyên tắc này cũng là tiêu chí để quyết định
một kết quả được đưa vào phân tích chính hay chỉ ghi ở phần bổ sung.

\subsection{Cơ sở lý thuyết của Transformer và cơ chế chú ý}

Transformer biểu diễn chuỗi token bằng các vector và dùng chú ý để xác định mức độ
liên quan giữa các vị trí \cite{Vaswani2017_attention}. Với ma trận truy vấn $Q$, khóa
$K$ và giá trị $V$, phép chú ý tích vô hướng có điều chỉnh được viết như sau:

\begin{equation}
  \operatorname{Attention}(Q,K,V)=
  \operatorname{softmax}\left(\frac{QK^{\mathsf{T}}}{\sqrt{d_k}}\right)V,
  \label{eq:scaled-attention}
\end{equation}

trong đó $d_k$ là số chiều của khóa. Phép chia cho $\sqrt{d_k}$ giúp giữ độ lớn của
tích vô hướng trong khoảng ổn định trước khi áp dụng softmax. Chú ý nhiều đầu thực
hiện nhiều phép chiếu song song rồi ghép kết quả, nhờ đó mô hình có thể học các kiểu
quan hệ khác nhau giữa các token. Trong mô hình ngôn ngữ, mỗi token được biểu diễn
theo ngữ cảnh của toàn bộ câu. Điều này đặc biệt hữu ích với tin tài chính, vì ý nghĩa
của một con số hoặc một động từ thường phụ thuộc vào doanh nghiệp, chủ thể và câu đứng
trước hoặc sau nó. FinBERT và PhoBERT là hai ví dụ về việc dùng mô hình nền thuộc họ
BERT cho miền tài chính hoặc tiếng Việt \cite{Devlin2019_bert,Yang2020_200608097v2,
Nguyen2020_200300744v3}.

Một bộ mã hóa Transformer thường gồm cơ chế chú ý nhiều đầu, mạng truyền thẳng theo
từng vị trí, kết nối tắt và chuẩn hóa. Vì cơ chế chú ý không tự chứa thông tin thứ tự,
mô hình cần biểu diễn vị trí hoặc một cách mã hóa tương đương. Các khối này giúp mô
hình vừa học quan hệ giữa những vị trí xa nhau, vừa duy trì đường truyền gradient ổn
định. PhoBERT cung cấp sẵn bộ tách từ và biểu diễn theo thiết kế của mô hình nền. Khi
dùng TFT, thứ tự ngày được giữ trong \texttt{TimeSeriesDataSet}; quy tắc khai báo biến
theo thời điểm có thể quan sát được trình bày trong tiểu mục về mô hình này.

\subsection{Cơ sở lý thuyết của PhoBERT}

PhoBERT là mô hình ngôn ngữ tiền huấn luyện dành cho tiếng Việt. Văn bản được tách từ
theo quy ước phù hợp rồi đưa qua bộ tách từ BPE để sinh chuỗi đơn vị từ con. Khi tinh
chỉnh cho bài toán cảm xúc, biểu diễn của bộ mã hóa được đưa qua tầng phân loại. Với ba
lớp, điểm logit $\mathbf{z}$ được chuyển thành xác suất:

\begin{equation}
  p_c=\frac{\exp(z_c)}{\sum_{k=1}^{3}\exp(z_k)},
  \qquad c\in\{\mathrm{NEG},\mathrm{NEU},\mathrm{POS}\}.
  \label{eq:sentiment-softmax}
\end{equation}

Việc chọn PhoBERT không loại bỏ yêu cầu thích nghi với miền tài chính. Tên công ty, mã
cổ phiếu, thuật ngữ chuyên ngành và cách viết phần trăm có thể khác với ngữ liệu tiền
huấn luyện. Vì vậy, tập tinh chỉnh phải ghi rõ quy tắc tách từ, giới hạn độ dài, cách
xử lý thực thể và loại trùng. Phần dữ liệu dùng để sinh đặc trưng cho dự báo cũng phải
được xử lý bằng điểm kiểm đã khóa theo quy tắc thời gian. Mô hình được kiểm tra trên tập
dữ liệu trong miền đã khóa, thay vì chỉ dựa vào độ chính xác trên tập hạt giống. Kết quả
phân loại cảm xúc cần được báo cáo tách biệt với kết quả dự báo để tránh đồng nhất hai
loại chất lượng khác nhau.

\subsection{Cơ sở lý thuyết của LSTM}

LSTM đọc chuỗi theo thứ tự thời gian và duy trì trạng thái ẩn $\mathbf{h}_t$ cùng trạng
thái ô nhớ $\mathbf{c}_t$ \cite{Hochreiter1997_lstm}. Với đầu vào $\mathbf{x}_t$, một
bước cập nhật được mô tả như sau:

\begin{align}
  \mathbf{i}_t &= \sigma(W_i\mathbf{x}_t+U_i\mathbf{h}_{t-1}+\mathbf{b}_i),\\
  \mathbf{f}_t &= \sigma(W_f\mathbf{x}_t+U_f\mathbf{h}_{t-1}+\mathbf{b}_f),\\
  \mathbf{o}_t &= \sigma(W_o\mathbf{x}_t+U_o\mathbf{h}_{t-1}+\mathbf{b}_o),\\
  \widetilde{\mathbf{c}}_t &= \tanh(W_c\mathbf{x}_t+U_c\mathbf{h}_{t-1}+\mathbf{b}_c),\\
  \mathbf{c}_t &= \mathbf{f}_t\odot\mathbf{c}_{t-1}+
                 \mathbf{i}_t\odot\widetilde{\mathbf{c}}_t,\\
  \mathbf{h}_t &= \mathbf{o}_t\odot\tanh(\mathbf{c}_t).
  \label{eq:lstm-update}
\end{align}

trong đó $\mathbf{i}_t$, $\mathbf{f}_t$ và $\mathbf{o}_t$ lần lượt là cổng vào, cổng
quên và cổng ra. Cổng quên quyết định phần thông tin cũ được giữ lại, cổng vào điều
chỉnh thông tin mới đi vào ô nhớ, còn cổng ra tạo trạng thái ẩn dùng ở bước tiếp theo.
Nhờ cơ chế này, LSTM có thể giữ những phụ thuộc kéo dài qua nhiều phiên mà không phải
truyền toàn bộ chuỗi vào một phép tổng hợp đơn giản.

LSTM phù hợp với mức sàn vì có thể mô hình hóa cửa sổ giá hoặc cửa sổ cảm xúc mà không
yêu cầu đồ thị liên mã. Trong hệ thống, hai nhánh LSTM dùng cùng số ngày quan sát, cùng
quy tắc đệm chuỗi hoặc loại mẫu và cùng điều kiện dừng để phép so sánh có ý nghĩa. Việc
thêm nhánh cảm xúc làm tăng số tham số, nên số chiều trạng thái và thời gian huấn luyện
cũng cần được ghi lại khi phân tích chênh lệch.

\subsection{Cơ sở lý thuyết của Bộ biến đổi hợp nhất theo thời gian}

TFT được đưa vào như mô hình mở rộng cho chuỗi nhiều biến \cite{Lim2021_tft}. Mô hình
kết hợp bộ mã hóa chuỗi với cơ chế chọn biến, các khối có cổng và chú ý. Nhờ đó, một
biến có thể có mức quan trọng khác nhau tùy thời điểm và trạng thái dự báo. Trong bài
toán này, động lượng, độ biến động, tỷ lệ cảm xúc và cờ có tin là những nhóm biến cần
được phân biệt.

Nếu được chạy, TFT phải được cấu hình cho bài toán phân loại. Đầu ra cuối là điểm logit
kích thước 3 và hàm mất mát là entropy chéo. Cấu hình hồi quy phân vị có sẵn trong thư
viện không được dùng thay cho mục tiêu này. Các biến cũng phải được khai báo rõ theo
trạng thái: biết trước, biến tĩnh hoặc chỉ quan sát được tại thời điểm dự báo. Tin tức
sau mốc cắt không được đưa vào nhóm biến đã biết tại ngày trước đó.

Tầm quan trọng của biến và thành phần chú ý chỉ được lưu như bằng chứng mô tả cách mô
hình sử dụng biến, chứ không phải ảnh hưởng nhân quả. Các nghiên cứu dùng Transformer
cho biến động và chuỗi đa phương thức cho thấy tiềm năng của hướng này, nhưng độ phức
tạp, số lượng tham số, khả năng quá khớp và khác biệt giữa các thị trường vẫn phải được
kiểm tra riêng \cite{Boyle2023_230503835v1,Omranpour2024_241210540v1}.

\subsection{Luồng xử lý tổng thể}

Luồng xử lý của hệ thống gồm các bước:

\begin{enumerate}[nosep]
  \item Khóa phạm vi mã, khoảng thời gian, mốc cắt, công thức nhãn và hạt giống;
  \item Thu thập tin tức, OHLCV và lịch giao dịch, đồng thời lưu nguồn gốc và khả năng truy xuất;
  \item Làm sạch, loại trùng, kiểm tra mốc thời gian và ánh xạ bài viết theo mã;
  \item Tạo nhánh cảm xúc và nhánh giá bằng các phép biến đổi độc lập;
  \item Căn chỉnh hai nhánh theo mã và ngày giao dịch, sau đó tạo nhãn ở thời điểm hợp lệ;
  \item Chạy thang mô hình từ đối chứng cơ sở đến LSTM hai nhánh, chỉ chạy TFT khi đủ điều kiện;
  \item Đánh giá bằng cuốn chiếu theo thời gian, loại bỏ đặc trưng và lấy mẫu lặp, rồi lưu bản kê kết quả.
\end{enumerate}

\ifShowDiagram
Hình~\ref{fig:kien-truc} minh họa các bước chính của quy trình xử lý. Các khối TFT và
chú ý chéo có cổng được đánh dấu là phần mở rộng, còn mô hình LSTM hai nhánh là cấu
hình cơ sở. Cách phân biệt này giúp người đọc nhận ra phần bắt buộc và phần chỉ được
thực hiện khi dữ liệu và thời gian cho phép.
\else
Các bước chính của quy trình xử lý được mô tả trong phần này; sơ đồ minh họa được lược
bỏ ở biến thể báo cáo không kèm hình.
\fi
\ifShowDiagram
\begin{figure}[H]
  \centering
  \begin{tikzpicture}[
      node distance=6mm and 10mm,
      font=\small,
      stage/.style={
        rectangle, rounded corners=2pt, draw=black!55, thick,
        text width=10.5cm, minimum height=9mm, align=center,
        fill=gray!6, inner sep=4pt},
      model/.style={
        rectangle, rounded corners=2pt, draw=blue!65, very thick,
        text width=10.5cm, minimum height=9mm, align=center,
        fill=blue!8, inner sep=4pt},
      floor/.style={
        rectangle, rounded corners=2pt, draw=green!55!black, very thick,
        text width=10.5cm, minimum height=9mm, align=center,
        fill=green!10, inner sep=4pt},
      ext/.style={
        rectangle, rounded corners=2pt, draw=violet!60, very thick, dashed,
        text width=10.5cm, minimum height=9mm, align=center,
        fill=violet!6, inner sep=4pt},
      arr/.style={-{Stealth[length=2.5mm]}, thick, black!65},
    ]
    \node[stage] (g1) {1. Thu thập dữ liệu\\ Giá OHLCV của 10 mã VN30 và tin tức Vietstock trong giai đoạn 2020--2025};
    \node[model, below=of g1] (g2) {2. Phân tích cảm xúc\\ PhoBERT được tinh chỉnh để sinh xác suất của 3 lớp: tích cực, trung tính và tiêu cực};
    \node[stage, below=of g2] (g3) {3. Đặc trưng giá\\ Lợi suất, lợi suất log, MA, khối lượng và biến động};
    \node[stage, below=of g3] (g4) {4. Căn chỉnh và gán nhãn\\ Gộp cảm xúc và giá theo mã và ngày, mốc cắt 15:00, nhãn UP, FLAT và DOWN};
    \node[floor, below=of g4] (g5) {5. Mô hình dự báo\\ Đối chứng cơ sở, sau đó là LSTM một nhánh và LSTM hai nhánh ở mức sàn tốt nghiệp};
    \node[ext, below=of g5] (g6) {Mở rộng có điều kiện: đặc trưng kỹ thuật và TFT};
    \node[stage, below=of g6] (g7) {6. Đánh giá\\ Cuốn chiếu theo thời gian + loại bỏ đặc trưng có/không cảm xúc, F1 vĩ mô, khoảng tin cậy lấy mẫu lặp};

    \draw[arr] (g1) -- (g2);
    \draw[arr] (g2) -- (g3);
    \draw[arr] (g3) -- (g4);
    \draw[arr] (g4) -- (g5);
    \draw[arr] (g5) -- (g6);
    \draw[arr] (g6) -- (g7);
  \end{tikzpicture}
  \caption{Kiến trúc tổng thể của quy trình dự báo xu hướng giá. Khối xanh dương thể
  hiện mô hình cảm xúc PhoBERT, khối xanh lá là mức sàn tốt nghiệp, khối nét đứt là
  phần mở rộng.}
  \label{fig:kien-truc}
\noindent\textit{Nguồn: Tác giả xây dựng.}
\end{figure}
\fi

\section{Thu thập và tiền xử lý dữ liệu}
\label{sec:pp-du-lieu}

\subsection{Phạm vi dữ liệu và nguồn thu thập}

Dữ liệu gồm hai bảng chính. Bảng tin tức lưu mã bản ghi, URL, nguồn, tiêu đề, nội
dung nếu có, mốc thời gian đăng hoặc cập nhật, thời điểm thu thập và mã cổ phiếu liên
quan. Bảng giá lưu mã, ngày giao dịch, OHLCV điều chỉnh và cờ chất lượng. Theo thiết
kế hiện hành, giai đoạn thực nghiệm được cố định từ năm 2020 đến hết năm 2025, với danh
sách 10 mã HOSE thuộc VN30 được ghi trong bản kê. Nếu phạm vi thực tế thay đổi theo độ
phủ dữ liệu, bản kê phải ghi mốc bắt đầu, mốc kết thúc và lý do thay đổi thay vì để
phạm vi chỉ xuất hiện trong phần mô tả.

Nguồn tin hiện được ưu tiên là Vietstock theo mã, vì cách tổ chức này giúp giảm mơ hồ
khi ánh xạ bài viết với cổ phiếu. CafeF, VnEconomy hoặc Stockbiz chỉ được bổ sung sau
khi kiểm tra độ phủ, nội dung, mốc thời gian và điều kiện sử dụng. Theo đề cương, các
nguồn này là nhóm nguồn dự kiến và có thể điều chỉnh khi khả năng truy cập hoặc cấu
trúc trang thay đổi. Dữ liệu giá được thu thập qua thư viện \texttt{vnstock} với nguồn
VCI theo cấu hình của đề tài. Một nguồn khác như \texttt{vnquant} chỉ dùng để đối chiếu
khi ngữ nghĩa của giá điều chỉnh, độ bao phủ và lược đồ đã được so sánh, và không tự
động được trộn vào tập chính. Việc tách nguồn chính và nguồn đối chiếu giúp phát hiện
sai lệch nhưng không làm tăng số quan sát bằng cách ghép các bản ghi chưa cùng quy ước.

\subsection{Kiểm tra chất lượng, nguồn gốc và khả năng truy xuất}

Mỗi bản ghi tin được gắn mã băm của văn bản đã chuẩn hóa cùng nguồn và URL. Bản ghi
trùng tiêu đề, nội dung hoặc có mốc thời gian gần nhau được đánh dấu để kiểm tra trước
khi loại. Việc loại trùng được thực hiện trước khi chia tập cảm xúc và trước khi tổng
hợp đặc trưng, vì một bản sao xuất hiện ở hai phần chia có thể làm kết quả tăng giả
tạo. Nếu một bài được gắn với nhiều mã hợp lệ, hệ thống giữ quan hệ nhiều mã và ghi rõ
cách mở rộng bản ghi. Một bài có thể được tính cho nhiều mã trong bảng liên kết, nhưng
không được xem là nhiều bài độc lập khi đánh giá độ không chắc chắn.

Nếu không xác định được mã hoặc thời điểm, bản ghi không được đưa vào bảng đặc trưng
chính. Những bản ghi bị loại vẫn được thống kê theo nguồn và lý do loại để người đọc
đánh giá độ bao phủ. Trường hợp mốc thời gian đăng và mốc thời gian cập nhật khác nhau,
cả hai trường đều phải được giữ nếu có; quy tắc chọn thời điểm sử dụng phải được khóa
trước khi tạo mẫu.

Các thống kê phải bao gồm số bài theo nguồn, mã và năm, tỷ lệ thiếu nội dung, tỷ lệ
thiếu mốc thời gian, số bản ghi trùng, độ phủ ngày có tin và số phiên giá hợp lệ. Cần
ghi thêm số bài gắn với một mã và nhiều mã, số ngày không có tin, số bản ghi bị chuyển
sang phiên kế tiếp và số bản ghi không thể xác minh. Bản kê cũng lưu thời điểm thu
thập, phiên bản mã nguồn, mã băm tệp dữ liệu và quy tắc làm sạch. Mục tiêu của bước
này là bảo đảm kho ngữ liệu có thể truy nguyên, thay vì chỉ tạo ra một tập dữ liệu lớn
nhưng khó kiểm tra.

\subsection{Dữ liệu giá}

Dữ liệu OHLCV được dùng theo giá điều chỉnh để tính lợi suất và hạn chế đứt gãy do chia
tách, cổ tức hoặc sự kiện doanh nghiệp. Với giá đóng cửa điều chỉnh
$C^{\mathrm{adj}}_{i,t}$ của mã $i$ tại ngày $t$, lợi suất mục tiêu được định nghĩa là:

\begin{equation}
  r_{i,t+1}=\frac{C^{\mathrm{adj}}_{i,t+1}}
  {C^{\mathrm{adj}}_{i,t}}-1.
  \label{eq:next-return}
\end{equation}

Các kiểm tra cơ bản gồm thứ tự thời gian, ngày trùng, giá không dương, khối lượng
không hợp lệ và phiên bị thiếu. Dữ liệu từ nguồn chính được đối chiếu với nguồn thay
thế trong phạm vi cần thiết để phát hiện ngày thiếu, bản ghi bất thường hoặc sai lệch
do thay đổi đơn vị. Các sự kiện chia tách, trả cổ tức bằng cổ phiếu hoặc điều chỉnh giá
được xem xét khi xây dựng chuỗi, vì dùng không nhất quán có thể tạo ra lợi suất bất
thường không phản ánh giao dịch thực tế.

Giá trị của ngày $t+1$ chỉ dùng để tạo nhãn, không dùng để tạo đặc trưng tại ngày $t$.
Những ngày đầu chuỗi chưa đủ dữ liệu cho trung bình động hoặc độ biến động được loại khỏi
mẫu hoặc xử lý theo quy tắc ghi trong bản kê. Mọi khác biệt giữa giá đóng cửa gốc
và giá điều chỉnh phải được giải thích trong báo cáo dữ liệu, không tự động coi một
chuỗi là sai chỉ vì hai nguồn có cách điều chỉnh khác nhau.

\subsection{Dữ liệu tin tức và ánh xạ theo phiên}

Mỗi bản tin giữ lại cả tiêu đề và nội dung để phục vụ kiểm tra, nhưng chính sách đầu
vào phải được khóa trước khi huấn luyện. Mức sàn ưu tiên tiêu đề vì phù hợp với tập
hạt giống công khai và giảm khác biệt về độ dài. Nội dung được lưu để kiểm tra, chỉ
được dùng trong phân tích độ nhạy khi đã được gán nhãn, tách theo thời gian và báo cáo
riêng. Khi một bài chỉ có tiêu đề, hệ thống không tự suy ra phần nội dung còn thiếu;
trường thiếu phải được lưu lại để phân tích độ bao phủ.

Tất cả mốc thời gian được chuyển về múi giờ Việt Nam. Với mỗi ngày giao dịch $t$, tin có
mốc thời gian không muộn hơn 15 giờ của ngày đó mới được xem là có thể quan sát tại mốc ra
quyết định. Tin sau 15 giờ được chuyển sang phiên giao dịch kế tiếp theo quy tắc bảo
thủ, nên không được dùng để tạo đặc trưng cho phiên đã đóng. Tin vào cuối tuần, ngày
nghỉ hoặc ngày không giao dịch cũng được gán sang phiên hợp lệ tiếp theo. Quy tắc này
có thể làm trễ một phần tín hiệu, nhưng giảm nguy cơ đánh giá lạc quan do dùng tin xuất
hiện sau mốc cắt.

Ánh xạ được thực hiện sau khi chuẩn hóa múi giờ và sắp xếp lịch giao dịch, không dựa
chỉ vào ngày in trong URL. Với tin có cả thời điểm đăng và thời điểm cập nhật, trường
được dùng cho dự báo phải được ghi trong bản kê. Mốc thời gian thiếu, mâu thuẫn hoặc
chỉ có ngày mà không có giờ không được gán một giờ mặc định rồi đưa thẳng vào tập
chính; bản ghi phải được đánh dấu để loại hoặc phân tích riêng.

\subsection{Tiền xử lý văn bản tiếng Việt}

Quy trình tiền xử lý văn bản gồm các bước sau:

\begin{enumerate}[nosep]
  \item Chuẩn hóa Unicode, khoảng trắng, ký hiệu phần trăm và cách viết số, nhưng giữ
    lại các ký hiệu có thể mang nghĩa tài chính;
  \item Tách từ bằng công cụ phù hợp với PhoBERT, giữ dấu tiếng Việt và các đơn vị từ tài
    chính có ý nghĩa như mã cổ phiếu, tên công ty và đơn vị tiền tệ;
  \item Nhận diện và che một số thực thể như tổ chức, địa danh hoặc mã cổ phiếu khi
    mục tiêu là học ngữ cảnh, đồng thời giữ bản gốc trong các trường nguồn gốc và
    truy xuất;
  \item Loại HTML, quảng cáo, phần điều hướng, bài rỗng và nội dung không thuộc văn bản
    tin theo quy tắc đã định, không tự tạo câu thay thế;
  \item Giới hạn độ dài theo điểm kiểm, dùng bộ tách từ BPE của PhoBERT và ghi lại
    chiều dài cùng số đơn vị từ bị cắt trong từng bản ghi.
\end{enumerate}

Số liệu, dấu phần trăm và ký hiệu tăng giảm không bị xóa máy móc vì chúng có thể mang
tín hiệu tài chính. Nếu sử dụng phép loại dấu câu hoặc che thực thể, quy tắc đó phải
được giữ nguyên giữa các phần chia. Mọi biến đổi nhìn thấy toàn bộ kho ngữ liệu, chẳng
hạn vốn từ mới hoặc thống kê chọn token, phải được loại khỏi thực nghiệm chính hoặc
chỉ được ước lượng trên phần huấn luyện.

\subsection{Tạo tập cảm xúc và gán nhãn}

Tập hạt giống gồm các tiêu đề tài chính đã được gán nhãn ba lớp, dùng để kiểm tra
nhanh quy trình. Tập dữ liệu trong miền được mở rộng theo một hướng dẫn gán nhãn cố
định. Người gán nhãn đọc phần nội dung được phép sử dụng, xác định chủ thể chính, sự
kiện và hướng tác động kỳ vọng đến giá trong ngắn hạn; nội dung đầy đủ chỉ được đọc
khi nó thuộc thời điểm và phạm vi được phép sử dụng. Các quy tắc cơ bản là:

\begin{itemize}[nosep]
  \item Tích cực nếu thông tin hàm ý triển vọng hoặc tác động có lợi rõ ràng;
  \item Tiêu cực nếu thông tin hàm ý rủi ro, suy giảm hoặc tác động bất lợi rõ ràng;
  \item Trung tính nếu chỉ thông báo, mô tả hoặc chưa đủ bằng chứng để xác định chiều.
\end{itemize}

Tin có hai chiều trái ngược được gán theo tác động nổi trội nếu hướng dẫn cho phép,
hoặc được đánh dấu để rà soát. Tin không xác định được chủ thể hoặc hướng tác động rõ
ràng được giữ lại với nhãn trung tính khi có đủ căn cứ; nếu không, mẫu được loại khỏi
tập huấn luyện và lý do được lưu lại.

Mỗi mẫu được gán một trong ba lớp NEGATIVE, NEUTRAL hoặc POSITIVE theo hướng dẫn đã
khóa. Báo cáo gán nhãn phải ghi số mẫu, tiêu chí, phiên bản hướng dẫn, số mẫu đã được
rà soát và phân bố lớp. Mô hình ngôn ngữ lớn chỉ tạo nhãn sơ bộ; một người duy nhất
rà soát toàn bộ nhãn trước khi huấn luyện chính thức, nên không tính được Cohen's hoặc
Fleiss' kappa do không có người gán nhãn thứ hai. Các mẫu dùng để đánh giá dự báo phải
được tách theo thời gian khỏi dữ liệu dùng để tinh chỉnh cảm xúc trong phần chia
tương ứng.

\section{Nhánh cảm xúc: tinh chỉnh PhoBERT}
\label{sec:pp-phobert}

\subsection{Thiết lập bài toán phân loại}

Mỗi bài viết hợp lệ là một mẫu có nhãn $y_j$ thuộc một trong ba lớp NEGATIVE, NEUTRAL
hoặc POSITIVE. Nhãn được hiểu theo hướng tác động tiềm năng của nội dung đối với doanh
nghiệp hoặc cổ phiếu liên quan, không phải chỉ theo cảm xúc chung của người viết. Mô
hình sử dụng điểm kiểm \texttt{vinai/phobert-base} làm bộ mã hóa, tiếp theo là tầng bỏ
ngẫu nhiên và tầng phân loại tuyến tính. Điểm logit được chuyển thành xác suất theo
Phương trình~\ref{eq:sentiment-softmax}. Tầng phân loại tuyến tính được chọn làm cấu
hình chính vì có ít tham số bổ sung và dễ tái lập. PhoBERT kết hợp CNN hoặc BiLSTM chỉ
được thêm vào khi dữ liệu và tiến độ cho phép, và chỉ được so sánh sau khi cấu hình cơ
sở đã được khóa.

PhoBERT được chọn theo ngôn ngữ của dữ liệu, còn FinBERT được dùng làm đối chiếu khái
niệm ở miền tiếng Anh \cite{Nguyen2020_200300744v3,Yang2020_200608097v2}. Mức sàn
không dịch toàn bộ tin sang tiếng Anh vì dịch máy có thể làm thay đổi sắc thái, tên
riêng và số liệu. Nếu có thí nghiệm dịch, nó phải là một đối chứng cơ sở riêng với
phiên bản công cụ và cấu hình được ghi lại.

\subsection{Chia dữ liệu và huấn luyện}

Xác thực chéo của PhoBERT dùng \texttt{StratifiedKFold} ngẫu nhiên với năm phần chia,
giữ tỷ lệ ba lớp trong mỗi tập kiểm thử ngoài. Phép đo này chỉ phản ánh hiệu năng
phân loại trong miền dữ liệu, chứ không phải đánh giá theo đúng thời điểm của nhánh dự
báo. Một phần của tập huấn luyện ngoài được rút ngẫu nhiên theo tỷ lệ lớp để chọn
điểm kiểm; tập kiểm thử ngoài không tham gia vào lựa chọn đó. Các bài có đầu vào trùng
khớp hoàn toàn bị loại trước khi chia tập để hạn chế việc mô hình ghi nhớ câu chữ và
tạo đánh giá lạc quan. Bản gần trùng hoặc các bài cùng nguồn chưa được phát hiện tự
động, nên vẫn là một giới hạn của giao thức.

Khi mô hình cảm xúc được dùng trong cuốn chiếu theo thời gian, có hai cách hợp lệ:
tinh chỉnh lại trong từng phần chia bằng nhãn của quá khứ, hoặc dùng một điểm kiểm
đóng băng được huấn luyện hoàn toàn trước giai đoạn kiểm thử. Cách được chọn phải
được ghi trong bản kê và áp dụng nhất quán cho tất cả mô hình so sánh. Không được
tinh chỉnh trên nhãn hoặc văn bản của tương lai rồi gọi đầu ra là ngoài mẫu. Nếu dữ
liệu gán nhãn không đủ cho việc tinh chỉnh lại ở mỗi phần chia, giới hạn này phải
được nêu cùng phạm vi diễn giải.

Hàm mất mát cơ bản là entropy chéo:

\begin{equation}
  \mathcal{L}_{\mathrm{CE}}=-\frac{1}{N}\sum_{j=1}^{N}
  \log p_{j,y_j}.
  \label{eq:sentiment-loss}
\end{equation}

Nếu lớp mất cân bằng, dùng trọng số $w_c$ được tính từ phần tập huấn luyện:

\begin{equation}
  \mathcal{L}_{\mathrm{WCE}}=-\frac{1}{N}\sum_{j=1}^{N}
  w_{y_j}\log p_{j,y_j}.
  \label{eq:weighted-sentiment-loss}
\end{equation}

Trọng số, cách lấy mẫu lại, hạt giống, kích thước lô, số vòng huấn luyện, bộ tách từ
và điểm kiểm đều được ghi trong bản kê. Điểm kiểm không được chọn dựa trên kết quả
của tập kiểm thử. Ngoài điểm số tổng hợp, cần giữ ma trận nhầm lẫn và các mẫu dự đoán
sai đại diện để phân tích những trường hợp phủ định, số liệu, tên riêng, tin có nhiều
sự kiện và bài trung tính.

\subsection{Đầu ra xác suất và tổng hợp theo mã, ngày}

Gọi $\mathcal{N}_{i,t}$ là tập tin hợp lệ của mã $i$ được ánh xạ vào ngày $t$, với
$n_{i,t}=|\mathcal{N}_{i,t}|$. Xác suất trung bình của lớp $c$ được tính bởi:

\begin{equation}
  \bar p_{i,t,c}=\frac{1}{n_{i,t}}
  \sum_{j\in\mathcal{N}_{i,t}}p_{i,j,t,c},
  \qquad n_{i,t}>0.
  \label{eq:daily-sentiment}
\end{equation}

Vector đặc trưng cảm xúc gồm $\bar p_{\mathrm{NEG}}$, $\bar p_{\mathrm{NEU}}$,
$\bar p_{\mathrm{POS}}$, chênh lệch $\bar p_{\mathrm{POS}}-\bar p_{\mathrm{NEG}}$,
số tin, độ lệch chuẩn của xác suất, tỷ lệ tin tích cực hoặc tiêu cực và biến nhị phân
\texttt{has\_news}. Có thể bổ sung các cửa sổ tổng hợp ngắn hạn như một ngày, ba ngày
hoặc năm ngày giao dịch gần nhất để kiểm tra khả năng tín hiệu kéo dài sau ngày công bố.

Ngày không có tin được gán vector xác suất theo tiên nghiệm trung tính, đặt số tin và
các tỷ lệ bằng không, đồng thời đặt \texttt{has\_news}=0 trong cấu hình mức sàn. Ngày
chỉ có tin trung tính vẫn có số tin dương và \texttt{has\_news}=1, nên hai trường hợp
được phân biệt. Ngoài cấu hình chính, có thể so sánh việc đặt vector bằng không hoặc
giữ giá trị gần nhất trong một giới hạn thời gian trên tập xác thực. Không được chọn
cách biểu diễn ngày không có tin dựa trên điểm của đoạn kiểm thử.

Các đặc trưng được tổng hợp sau mốc cắt, và số tin hay xác suất của bài đăng sau mốc
cắt không được dùng để mô tả ngày đã kết thúc. Nếu một bài gắn với nhiều mã, cách
nhân bản đặc trưng phải được ghi rõ để không đếm bài như nhiều quan sát độc lập khi
phân tích độ không chắc chắn. Mức độ phủ tin, số lượng bài và tỷ lệ bản ghi bị loại
cũng được lưu để khi diễn giải một ngày ít tín hiệu, người đọc có thể phân biệt giữa
``không có tin'' và ``có tin nhưng mô hình không chắc chắn''.

\section{Nhánh giá và ghép đặc trưng}
\label{sec:pp-ghep}

\subsection{Đặc trưng kỹ thuật}

Từ dữ liệu OHLCV điều chỉnh, hệ thống tạo các nhóm đặc trưng sau:

\begin{itemize}[nosep]
  \item Lợi suất ngày và lợi suất log;
  \item Tỷ lệ giữa giá đóng cửa với trung bình động năm phiên;
  \item Thay đổi khối lượng trong một phiên;
  \item Độ lệch chuẩn lợi suất trong cửa sổ năm phiên để mô tả độ biến động;
  \item Số lượng tin, xác suất cảm xúc tổng hợp, tỷ lệ lớp và biến chỉ báo có tin.
\end{itemize}

Các đặc trưng được tính từ dữ liệu đến ngày quan sát $t$. Để giảm phụ thuộc vào mức
giá tuyệt đối, các đại lượng có đơn vị giá được chuyển thành tỷ lệ hoặc sai khác tương
đối. Mỗi đặc trưng chỉ nhìn về quá khứ, không dùng giá của phiên kế tiếp. Các hàng
đầu chuỗi chưa đủ lịch sử được giữ ở dạng khuyết và bị loại khi tạo cửa sổ. Bộ chuẩn
hóa được ước lượng một lần trên phần huấn luyện rồi áp dụng nguyên trạng cho xác thực
và kiểm thử.

\subsection{Căn chỉnh hai nguồn dữ liệu}

Mỗi dòng bảng đặc trưng có khóa $(i,t)$ gồm vector giá $\mathbf{x}^{p}_{i,t}$, vector
cảm xúc $\mathbf{x}^{s}_{i,t}$, cờ chất lượng, thông tin mã và nhãn mục tiêu. Với
cửa sổ quan sát dài $L$, đầu vào của mô hình tại ngày $t$ là:

\begin{equation}
  X^{p}_{i,t}=[\mathbf{x}^{p}_{i,t-L+1},\ldots,\mathbf{x}^{p}_{i,t}],
  \qquad
  X^{s}_{i,t}=[\mathbf{x}^{s}_{i,t-L+1},\ldots,\mathbf{x}^{s}_{i,t}].
  \label{eq:input-windows}
\end{equation}

Các mã được chia theo ngày để bảo đảm mọi mẫu trong một phần chia dùng cùng mốc lịch.
Nếu một ngày không có tin, vector cảm xúc vẫn được tạo theo quy tắc ở trên. Không
được loại ngày một cách có chọn lọc vì việc đó có thể làm sai lệch phép so sánh giữa
giai đoạn nhiều tin và giai đoạn ít tin. Những ngày thiếu giá hoặc thiếu lịch sử đủ
dài phải được ghi rõ lý do loại; không được loại riêng các ngày khó dự báo.

\subsection{Gán nhãn xu hướng phiên kế tiếp}

Lợi suất mục tiêu được tính theo Phương trình~\ref{eq:next-return}. Với cửa sổ huấn
luyện $w$, hai phân vị $q^-_w$ và $q^+_w$ được chọn từ lợi suất chỉ trong phần huấn
luyện. Quy tắc phân lớp là:

\begin{equation}
 y_{i,t+1}=\begin{cases}
 \mathrm{DOWN}, & r_{i,t+1}<q^-_w,\\
 \mathrm{FLAT}, & q^-_w\leq r_{i,t+1}\leq q^+_w,\\
 \mathrm{UP}, & r_{i,t+1}>q^+_w.
 \end{cases}
 \label{eq:trend-label}
\end{equation}

Trong giao thức chính, $q^-_w$ và $q^+_w$ được lấy từ hai phân vị thấp và cao đã khóa
trong cấu hình. Nếu dùng ngưỡng cố định, ngưỡng đó phải được ghi rõ và không được
điều chỉnh dựa trên kết quả kiểm thử. Việc chọn ngưỡng cũng phải cân bằng hai yêu
cầu: lớp đi ngang cần phản ánh các biến động nhỏ chưa có hướng rõ rệt, nhưng phân bố
ba lớp không nên lệch đến mức mô hình chỉ cần dự đoán một lớp phổ biến là đạt điểm
cao. Ước lượng ngưỡng trên tập huấn luyện giúp phân phối lợi suất tương lai không
quyết định trước tỷ lệ lớp của giai đoạn quá khứ. Phân bố nhãn được báo cáo theo mã
và từng giai đoạn để người đọc đánh giá độ khó của bài toán.

\subsection{Hợp nhất bằng LSTM hai nhánh}

Nhánh giá và nhánh cảm xúc sử dụng hai bộ mã hóa LSTM riêng:

\begin{equation}
  \mathbf{h}^{p}_{i,t}=\operatorname{LSTM}_{p}(X^{p}_{i,t}),
  \qquad
  \mathbf{h}^{s}_{i,t}=\operatorname{LSTM}_{s}(X^{s}_{i,t}).
  \label{eq:two-branch-encoders}
\end{equation}

Hai trạng thái cuối được ghép bằng phép nối và đưa qua tầng phân loại hợp nhất:

\begin{equation}
  \mathbf{z}_{i,t}=\left[\mathbf{h}^{p}_{i,t};\mathbf{h}^{s}_{i,t}\right],
  \qquad
  \hat{\mathbf{y}}_{i,t+1}=\operatorname{softmax}
  \left(W_2\,\phi(W_1\mathbf{z}_{i,t}+\mathbf{b}_1)+\mathbf{b}_2\right),
  \label{eq:fusion-classifier}
\end{equation}

trong đó $\phi$ là hàm kích hoạt và $\hat{\mathbf{y}}$ là xác suất của ba lớp. Mô
hình chỉ giá dùng nhánh giá và tầng phân loại tương ứng nhưng không có
$\mathbf{h}^{s}$. Khi so sánh, số chiều ẩn, cửa sổ, hạt giống, ngân sách huấn luyện,
tiêu chí dừng sớm và cách ước lượng tiền xử lý phải gần như nhau. Chênh lệch tham số
do nhánh cảm xúc phải được báo cáo để người đọc đánh giá đúng mức độ khác nhau giữa
hai cấu hình. Nếu hai cấu hình còn khác biệt ngoài sự có mặt của cảm xúc, khác biệt
đó phải được xem là một yếu tố gây nhiễu trong phép loại bỏ đặc trưng.

\section{Thang mô hình dự báo}
\label{sec:pp-mo-hinh}

Đồ án được tổ chức theo nguyên tắc thang tầng. Mỗi tầng trả lời một câu hỏi kiểm
chứng riêng và chỉ được mở rộng sau khi tầng trước có dữ liệu hợp lệ.

\subsection{Tầng đối chứng cơ sở}

Tầng đầu gồm đối chứng lớp phổ biến nhất và đối chứng ngẫu nhiên. Đối chứng lớp phổ
biến nhất cho biết mức điểm đạt được khi luôn dự đoán lớp phổ biến nhất trong phần
huấn luyện. Đối chứng ngẫu nhiên được sinh theo quy tắc đã khóa, có hạt giống và
chạy trên cùng tập kiểm thử. Khi cần, có thể báo cáo thêm đối chứng ngẫu nhiên đều
để phân biệt ảnh hưởng của phân bố lớp với ảnh hưởng của mô hình. Hai mốc này không
phải mô hình học quy luật từ dữ liệu, nhưng cần thiết để kiểm tra một kiến trúc có
thực sự vượt qua dự đoán đơn giản hay chỉ tận dụng sự mất cân bằng lớp.

\subsection{Tầng mô hình cổ điển chỉ giá}

Tầng này dùng một bộ phân loại phi tuần tự nhận vector đặc trưng giá của phiên quan
sát cuối cùng trong cửa sổ, nhằm tạo mốc so sánh cho biết dữ liệu có tín hiệu đủ rõ
trước khi dùng mô hình sâu hay không. Cấu hình đã thực hiện là hồi quy logistic đa
lớp. Rừng ngẫu nhiên và XGBoost được nêu trong đề cương nhưng không được chạy: chúng
trả lời cùng một câu hỏi với hồi quy logistic, trong khi câu hỏi trung tâm của đề
tài là phần đóng góp của cảm xúc, nên ngân sách thực nghiệm được dành cho phép đối
chứng đó.

Tầng cổ điển được chạy ở cả hai dạng, chỉ giá và giá kèm cảm xúc, trên đúng các dòng
mà mô hình tuần tự sử dụng, để hai họ mô hình được chấm trên cùng tập mẫu. Dạng chỉ
giá không nhận bất kỳ cột cảm xúc nào, kể cả số tin, cờ có tin hoặc siêu dữ liệu phát
sinh từ văn bản.

\subsection{Tầng LSTM chỉ giá}

LSTM chỉ giá nhận $X^p$ trong một cửa sổ cố định và dự báo ba lớp. Đây là đối chứng
cơ sở theo hướng học sâu, dùng để tách câu hỏi ``LSTM có khai thác được lịch sử giá
hay không'' khỏi câu hỏi ``cảm xúc có bổ sung thông tin hay không''. Siêu tham số
được chọn trên tập xác thực và giữ nguyên khi chạy loại bỏ đặc trưng. Đầu vào chỉ
gồm những đặc trưng có thể quan sát đến ngày $t$, còn nhãn được tạo từ lợi suất của
phiên kế tiếp.

\subsection{Tầng LSTM hai nhánh có cảm xúc}

Đây là mức sàn của đề tài. Nhánh cảm xúc nhận chuỗi đặc trưng tổng hợp theo ngày,
nhánh giá nhận chuỗi chỉ báo kỹ thuật; hai trạng thái được ghép bằng phép nối rồi đưa
qua tầng phân loại. Nhánh cảm xúc không nhận giá của phiên tương lai và nhánh giá
không nhận siêu dữ liệu của bài viết ngoài các đặc trưng đã được quy định. Thí
nghiệm loại bỏ đặc trưng chạy hai cấu hình:

\begin{enumerate}[nosep]
  \item Chỉ giá: chỉ sử dụng $X^p$;
  \item Giá kết hợp cảm xúc: sử dụng cả $X^p$ và $X^s$.
\end{enumerate}

Hai cấu hình phải dùng cùng ngày huấn luyện, xác thực và kiểm thử, cùng cách ước
lượng các phép tiền xử lý, cùng hạt giống hoặc tập hạt giống, cùng quy tắc nhãn và
cùng chỉ số. Nếu cảm xúc không cải thiện chỉ số, kết quả vẫn được giữ lại và phân
tích theo độ ổn định, thay vì thay đổi giao thức để tìm một bảng điểm có lợi.

\subsection{Tầng TFT mở rộng}

TFT chỉ được chạy sau khi các đối chứng cơ sở và LSTM hai nhánh đã có số liệu và mốc
cắt không còn lỗi. Mô hình nhận các biến thời gian, biến tĩnh hoặc biến quan sát theo
cấu hình rõ ràng, dùng kích thước đầu ra bằng 3 và hàm mất mát phân loại. Các biến
không được khai báo chỉ vì chúng có sẵn trong một bảng dữ liệu, mà phải được kiểm tra
theo thời điểm có thể quan sát. Kết quả TFT chỉ được so sánh với các tầng trước khi
dùng cùng các đoạn kiểm thử và quy tắc tạo nhãn. Trong đồ án này TFT được cài đặt
tối giản bằng PyTorch thuần (\texttt{TemporalFusionClassifier}): nhúng từng biến,
mạng chọn biến theo thời điểm, bộ mã hóa LSTM một tầng, chú ý đa đầu với cổng phần dư
và đầu phân loại GRN; cấu hình hồi quy phân vị của thư viện không được dùng cho mục
tiêu ba lớp.

\subsection{Huấn luyện và kiểm soát độ phức tạp}

Mô hình được huấn luyện bằng entropy chéo cho ba lớp, có dừng sớm theo F1 vĩ mô trên
tập xác thực. Tầng bỏ ngẫu nhiên, suy giảm trọng số và cắt biên độ gradient có thể
được dùng để hạn chế quá khớp, nhưng lựa chọn và số lần thử phải được ghi lại. Nếu
lớp mất cân bằng, trọng số hàm mất mát hoặc cách lấy mẫu lại chỉ được tính từ tập
huấn luyện. Mỗi mô hình lưu điểm kiểm tốt nhất theo chỉ số đã chọn, cùng hạt giống,
cửa sổ, số tham số và thời gian huấn luyện. Tập kiểm thử không được dùng để điều
chỉnh kiến trúc sau mỗi lần chạy.

\section{Giao thức đánh giá chống rò rỉ}
\label{sec:pp-danh-gia}

Đánh giá là phần trung tâm của phương pháp. Một phép chia tập theo thời gian chưa đủ
bảo đảm kết quả sạch nếu bộ chuẩn hóa, ngưỡng, phép chọn đặc trưng hoặc điểm kiểm đã
nhìn thấy dữ liệu tương lai. Giao thức vì vậy dùng cuốn chiếu theo thời gian với các
đoạn kiểm thử tách biệt, đồng thời kiểm tra toàn bộ các phép biến đổi.

\subsection{Cuốn chiếu theo thời gian}

Mọi mã được chia theo cùng một mốc ngày. Một cửa sổ cuốn chiếu theo thời gian gồm phần
huấn luyện ở quá khứ, phần xác thực ở đoạn kế tiếp và phần kiểm thử ở đoạn sau đó. Có
thể dùng cửa sổ mở rộng để tăng dần dữ liệu huấn luyện hoặc cửa sổ trượt để giới hạn
lịch sử; lựa chọn, độ dài các đoạn và bước dịch cửa sổ phải được ghi trong bản kê.
Mỗi cửa sổ phải chỉ rõ ngày cuối của tập huấn luyện, ngày bắt đầu và kết thúc của
xác thực, ngày bắt đầu và kết thúc của kiểm thử, cùng số mẫu thực tế sau khi tạo cửa
sổ lịch sử.

Cấu hình đã thực hiện dùng năm cửa sổ mở rộng, mỗi cửa sổ có 60 ngày xác thực và 60
ngày kiểm thử, các đoạn kiểm thử không chồng lấn nhau. \textbf{Không có tập giữ lại
cuối tách riêng}, vì cũng \textbf{không có bước tìm kiếm siêu tham số}: cấu hình mô
hình được cố định trước khi chạy và không thay đổi sau khi xem kết quả. Trong mỗi cửa
sổ, tập xác thực chỉ dùng để chọn điểm dừng huấn luyện, còn đoạn kiểm thử được chấm
đúng một lần. Vì không có cấu hình nào được tinh chỉnh trên đoạn kiểm thử, một tập
giữ lại thứ hai không bổ sung thêm bảo đảm nào; nêu rõ điều này phản ánh đúng thiết
kế hơn là mô tả một tập giữ lại không được dùng trong kết quả.

Với cửa sổ mở rộng, không trộn mẫu của một ngày giữa tập huấn luyện và kiểm thử. Các
mã cùng ngày phải nằm trong cùng một phần chia để tránh việc một mã tiết lộ trạng
thái thị trường cho mã khác qua cách chia lệch. Mẫu được tạo từ cửa sổ lịch sử cũng
phải được cắt theo ngày mục tiêu, không chỉ theo ngày bắt đầu cửa sổ. Nếu một mẫu
dùng lịch sử có phần nằm ở giai đoạn trước nhưng nhãn mục tiêu nằm sau mốc chia, mẫu
đó thuộc phần chia của nhãn.

Trong cấu hình hiện tại, các đoạn chia theo ngày mục tiêu được đặt liền nhau và không
chèn khoảng cấm. Với nhãn phiên kế tiếp, giá dùng để tạo đặc trưng tại ngày quan sát
đã được biết trước khi bắt đầu phiên mục tiêu; nhãn của đoạn sau không được dùng để
huấn luyện đoạn trước. Nếu mở rộng chân trời dự báo hoặc dùng đặc trưng công bố trễ,
phải bổ sung khoảng cấm và kiểm tra lại giao thức.

\subsection{Kiểm soát thông tin tại điểm dự báo và các phép biến đổi}

Danh mục kiểm tra rò rỉ gồm các hạng mục sau:

\begin{itemize}[nosep]
  \item Mốc thời gian tin và mốc cắt 15 giờ;
  \item Ánh xạ cuối tuần, ngày nghỉ và tin sau giờ;
  \item Giá điều chỉnh và công thức lợi suất mục tiêu;
  \item Bộ chuẩn hóa, bộ điền khuyết và phép chọn đặc trưng chỉ ước lượng trên tập huấn luyện;
  \item Ngưỡng $q^-_w$, $q^+_w$ chỉ tính từ lợi suất tập huấn luyện;
  \item Nhãn cảm xúc dùng để tinh chỉnh và điểm kiểm của từng phần chia;
  \item Loại trùng không để cùng một bài hoặc bản sao gần giống xuất hiện ở hai phần chia;
  \item Lựa chọn siêu tham số chỉ dùng tập huấn luyện và xác thực, không dùng đoạn kiểm thử.
\end{itemize}

Mỗi phần chia cần lưu một số ví dụ kiểm tra thủ công, gồm tin trước mốc cắt, tin sau
mốc cắt, tin cuối tuần và ngày không có tin. Cần kiểm tra cả trường hợp bài gắn với
nhiều mã, bài gần trùng giữa hai nguồn, ngày đầu chuỗi thiếu chỉ báo và mẫu có nhãn
ở phiên kế tiếp. Nếu bản ghi không thể xác minh, bản ghi đó bị loại khỏi phân tích
chính và tỷ lệ loại phải được báo cáo.

\subsection{Các chỉ số đánh giá}

Từ ma trận nhầm lẫn của ba lớp, độ chính xác dương tính, độ bao phủ và F1 được tính
riêng cho từng lớp. F1 vĩ mô là trung bình số học của F1 ba lớp:

\begin{equation}
  \operatorname{MacroF1}=\frac{1}{3}\sum_{c=1}^{3}F1_c.
  \label{eq:macro-f1}
\end{equation}

Độ chính xác cân bằng là trung bình độ bao phủ của ba lớp. AUC-ROC một-chống-còn-lại
lần lượt xem mỗi lớp là tích cực và hai lớp còn lại là tiêu cực, rồi lấy trung bình
các giá trị AUC. Độ chính xác hướng được báo cáo như tỷ lệ dự đoán đúng lớp, nhưng
không được dùng một mình khi phân bố lớp mất cân bằng. Các chỉ số được tính trên dự
đoán ngoài mẫu của từng phần chia, sau đó tổng hợp theo quy tắc đã khóa.

\subsection{Loại bỏ đặc trưng và khoảng tin cậy}

Phép loại bỏ đặc trưng chính so sánh cấu hình chỉ giá với cấu hình giá kết hợp cảm
xúc. Hai mô hình dùng cùng ngày kiểm thử để tạo chênh lệch theo cặp:

\begin{equation}
  \Delta_m= m(\mathrm{price+sentiment})-m(\mathrm{price\text{-}only}),
  \label{eq:ablation-delta}
\end{equation}

với $m$ là F1 vĩ mô, độ chính xác cân bằng hoặc AUC-ROC một-chống-còn-lại. Khoảng
tin cậy được ước lượng bằng hai cách lấy mẫu lặp, cả hai đều được báo cáo. Cách thứ
nhất lấy mẫu lặp trên các cửa sổ; cách thứ hai lấy mẫu lặp trên các \emph{ngày} kiểm
thử, trong đó mọi mã của cùng một ngày được lấy cùng nhau để giữ tương quan chéo
trong một phiên. Cách thứ hai được ưu tiên khi số cửa sổ nhỏ, vì khoảng tin cậy dựng
từ chỉ năm khối sẽ rất thô và có thể loại trừ 0 một cách ngẫu nhiên. Số lần lấy mẫu
lặp và hạt giống phải được ghi lại. Nếu khoảng tin cậy của $\Delta_m$ chứa 0, kết
luận phù hợp là chưa có bằng chứng rõ rằng cảm xúc cải thiện chỉ số đó trong giao
thức hiện tại, chứ không phải kết luận cảm xúc không hữu ích trong mọi bối cảnh.

Phương trình~\eqref{eq:ablation-delta} tự nó chưa tách được nguồn gốc của chênh lệch,
vì mô hình hai nhánh khác mô hình chỉ giá ở cả kiến trúc lẫn thông tin đầu vào. Đối
chứng bổ sung giữ nguyên kiến trúc hai nhánh, bài tin, thời điểm và khối lượng tin,
rồi thay vector xác suất bằng một phân phối tiên nghiệm trung tính hằng số. Khi đó
\begin{align}
  \substack{\text{hiệu ứng kiến trúc, sự hiện diện}\\\text{và khối lượng tin}} &= m(\text{hai nhánh, tiên nghiệm trung tính})\notag\\
  &\quad - m(\text{chỉ giá}),\\
  \text{đóng góp thông tin phân cực} &= m(\text{hai nhánh, cảm xúc thật})\notag\\
  &\quad - m(\text{hai nhánh, tiên nghiệm trung tính}),
\end{align}
và tổng của hai thành phần bằng đúng $\Delta_m$. Điều kiện để phép tách có hiệu lực
là hai lần chạy dùng cùng cấu hình, cùng cửa sổ, cùng hạt giống và cho dự đoán trùng
khớp ở cấu hình chỉ giá.

\subsection{Phân tích độ bền và điều kiện phát huy tín hiệu}

Sau khi bảng mức sàn đã ổn định, có thể phân tầng kết quả theo ngày có nhiều tin
hoặc ít tin, theo độ biến động thấp hoặc cao, theo trạng thái tăng hoặc giảm, và
theo mã cổ phiếu. Các phân tích này chỉ được xem là bổ sung vì làm tăng số phép so
sánh. Số lần thử và tiêu chí phân tầng cần được ghi rõ để tránh chọn riêng một nhóm
có kết quả tốt. Kết quả theo mã có thể có phương sai lớn khi một số mã có ít tin, vì
vậy cần báo cáo cả số quan sát của từng nhóm.

\subsection{Giải thích mô hình}

Với LSTM hai nhánh, giải thích cơ bản dựa trên phép loại bỏ đặc trưng hoặc che khuất:
lần lượt che nhóm cảm xúc, số tin hoặc một nhóm chỉ báo rồi đo thay đổi chỉ số. Với
TFT, mức độ quan trọng của biến và thành phần chú ý được lưu như thông tin bổ sung.
Không phương pháp nào trong số này cho phép coi biến cảm xúc là nguyên nhân của lợi
suất. Vì vậy, diễn giải nên dùng các cụm như ``mô hình đặt trọng số cao hơn'' hoặc
``chỉ số thay đổi khi che biến'', không dùng cụm ``cảm xúc gây ra biến động''.

\subsection{Tái lập và báo cáo}

Mỗi lần chạy cần lưu hạt giống, phiên bản Python và thư viện, điểm kiểm, danh sách
mã, mã băm dữ liệu, thời điểm thu thập, cấu hình mô hình, phép chia tập, chỉ số và
số lần thử. Báo cáo phải phân biệt ba trạng thái: đã đo, chưa chạy và không áp dụng.
Ô chưa có kết quả không được điền số tham chiếu từ công trình khác. Cùng với bảng
điểm, cần báo cáo độ phủ tin, số bản ghi bị loại, phân bố lớp và các giới hạn về
nguồn dữ liệu.

Bảng~\ref{tab:cau-hinh} tóm tắt các quyết định thiết kế chính của quy trình.

\begin{table}[H]
  \centering
  \caption{Tóm tắt các quyết định thiết kế chính.}
  \label{tab:cau-hinh}
  \begin{tabular}{p{4cm} p{9cm}}
    \toprule
    Thành phần & Lựa chọn \\
    \midrule
    Mô hình cảm xúc & PhoBERT-base tinh chỉnh, ba lớp và đầu ra xác suất \\
    Mô hình chuỗi thời gian & LSTM hai nhánh là mức sàn, TFT là tầng mở rộng \\
    Hợp nhất & Nối hai trạng thái cuối rồi phân loại \\
    Phạm vi & 10 mã VN30 trên HOSE, dữ liệu từ năm 2020 đến hết năm 2025 \\
    Mốc cắt & 15 giờ theo múi giờ Việt Nam, tin sau giờ đóng cửa chuyển sang phiên kế tiếp \\
    Nhãn & Lợi suất phiên kế tiếp, ngưỡng chỉ ước lượng trên tập huấn luyện \\
    Đánh giá & Cuốn chiếu theo thời gian, năm đoạn kiểm thử tách biệt, F1 vĩ mô là chỉ số chính \\
    \bottomrule
  \end{tabular}
\noindent\textit{Nguồn: Tác giả xây dựng.}
\end{table}
